% answers_short.tex —— 1977 年·安徽卷（理科） · 简短答案速查
% 由 tools/make_short_answers.py 自 answers.tex 的 \gksolhead/\gkanswer 提取，
% 并与 1977/1977_数学真题_安徽理_答案要点.md 逐题比对后生成，勿手改（改源 MD 后重跑）。

\begin{gknotice}
\textbf{使用说明}：本节只给最终答案便于快速核对，完整推导见解析版。答案为编写者独立推导，\textbf{非官方标准答案}。
\end{gknotice}

\gkpart{题目}
\begin{itemize}[leftmargin=2.6em]
  \item[\textbf{1.}] (1) $0$；(2) $x>2$；(3) $\dfrac{\sqrt6-\sqrt2}{4}$；(4) $0\leqslant x\leqslant1$；(5) $\sqrt2\left(\cos\dfrac{3\pi}{4}+i\sin\dfrac{3\pi}{4}\right)$；(6) 按题面为 $-\sin A\cos A=-\dfrac12\sin2A$；源解答按 $\sec A$ 口径得 $-1$（见解析）；(7) $9.42\ \mathrm{cm^2}$
  \item[\textbf{2.}] (1) $y=(80-2x)x=-2x^2+80x$（$0<x<40$）；(2) 顶点坐标 $(20,800)$，对称轴方程 $x=20$；(3) 槽深为 $20$ 厘米时截面积最大，最大面积为 $800\ \mathrm{cm^2}$
  \item[\textbf{3.}] (1) $T_{r+1}=(-1)^r\mathrm{C}_n^r\left(\dfrac23\right)^r x^{\frac{5n-8r}{15}}$；(2) 见解析版
  \item[\textbf{4.}] (1) $\lg x-1$（$x\geqslant10$）；$1-\lg x$（$1<x\leqslant10$）；(2) $S_n=\dfrac{n(n-1)}{2}+(1-2^n)\lg3$
  \item[\textbf{5.}] 三个内角为 $45^\circ$、$60^\circ$、$75^\circ$。
  \item[\textbf{6.}] (1) $y=x-6$；(2) 交点为 $(4,-2)$ 与 $(9,3)$；(3) 见解析版
  \item[\textbf{7.}] 轨迹为椭圆 $\dfrac{x'^2}{16}+\dfrac{y'^2}{9}=1$；略图见解析版。
  \item[\textbf{8.}] (1) $243$ 个；(2) 见解析版
\end{itemize}
